\section{Conclusions}\label{sec:conclusions}

A server that reads a connection's first bytes and decides its protocol on one listening socket can
match the same server with one port per protocol. For \oneport, under load of one protocol at a
time, one-port mode was equivalent to dedicated mode within \MarginLow{} to \MarginHigh{} in
\CostPassed{} of the \CostHolm{} tested cost cells. The claim of no measurable cost holds for \CostClaimCount{}
configurations: HTTP/1.1, TLS and MQTT on both Linux backends, and HTTP/1.1 and TLS on Windows. The
TLS cells have little power to show the cost of detection, since the handshake dominates. The
claim does not hold for h2c on L, where equivalence of the time to first byte at a fixed load was not shown. On epoll
the analysis reports a measured cost there. On W the claim is open for h2c and MQTT.

The detector decided all \HardcaseRuns{} runs of the frozen hard cases as specified, with no timed
event out of bounds. By B3's test, nginx, HAProxy and Envoy in both cases, and hyper-util in the
silent case, held at least \FootprintMargin{} times the server's counted footprint per pending
connection, as configured. Their medians of the difference range from
\BThreeSilentHyperUtilIouringDMedian{} to \BThreePartialHelloEnvoyEpollDMedian{} bytes, with the
intervals and exclusions of Section~\ref{sec:res-robust}. sslh-ev held less than the server in
every cell, as did Netty in the partial-ClientHello case. Replay served more connections per second than peek in every cell but h2c
on IOCP. In-process dispatch used less CPU than relay for HTTP/1.1. The server's relay served more
connections per second than four of the five proxies, and sslh-ev could not be tested.

The runs also exposed a limit of the server's io\_uring backend: beside a busy keep-alive background,
its churn starved in both modes. A larger ring of provided buffers is the first follow-up, labelled
exploratory, and the cause of h2c's latency at a fixed load is the second.
